\section{Introduction}

Experimental systems need more than an emulator. A run also has an input
schedule, a stopping policy, a choice of observations, a way to deliver those
observations, and a policy for retaining results. When a machine backend
implements all of these concerns together, downstream tools inherit its
telemetry conventions and storage assumptions. Changing what is observed can
then require changes to the execution engine, while comparing two machines
can require translating between unrelated result formats.

This coupling becomes more costly across heterogeneous virtual machines.
Different machines expose different artifacts, configuration fields, input
coordinates, state, and native evidence. A shared experiment still needs to
control runs consistently, but a common interface must not erase these
machine-specific semantics or imply that all machines expose the same
observations. It must make the common parts explicit and leave the rest with
the bundle that understands the machine.

GlassVM provides that separation. A machine bundle declares its identity and
contract, accepts a caller-owned execution request, and supplies an emulator
backend. Preparation validates artifact, configuration, inputs, controls,
and requested observations before execution. The resulting observation
contract selects typed evidence channels and capability dependencies.
Execution emits the selected evidence synchronously to caller-owned sinks.
Normalizers derive requested capability outputs with bounded streaming state,
and receipts account for evidence fulfillment, sink delivery, recorder
finalization, and publication as distinct outcomes.

The design keeps downstream policy outside the emulator. A caller may consume
emissions directly, persist them through the reference recorder, or do both.
GlassVM does not assemble a canonical whole-run trace or require every
consumer to retain a run. The recorder is one explicit consumer with its own
bounded buffering and publication protocol; its output is not the only form
of evidence available to an experiment.

We study this contract across CHIP-8, PICO-8, and TIC-80. CHIP-8 supplies a
bytecode virtual-machine baseline. PICO-8 and TIC-80 are independent
compatibility implementations with bounded Lua runtimes and documented
subsets. They provide materially different execution and observation paths,
while their fidelity limits remain visible in the corresponding machine
contracts.

This paper makes four contributions:

\begin{itemize}[leftmargin=*]
  \item A machine-neutral execution and preparation contract that separates
        bundle-specific semantics from experiment orchestration.
  \item An explicit observation model in which evidence requirements are
        negotiated before execution and fulfilled through typed bounded
        emissions.
  \item A reference durable recorder with separate evidence, recorder, and
        publication accounting.
  \item A three-machine evaluation spanning CHIP-8, PICO-8, and TIC-80,
        including clean-room package consumption, Python sink delivery,
        semantic equivalence, bounded-memory measurements, and a small
        downstream consumer.
\end{itemize}

The empirical magnitudes and exact source coordinates remain unresolved in
this draft. They will be supplied only from sealed evidence and source
receipts; no intermediate experiment summary is treated as final evidence.
